\documentclass[12pt, letterpaper]{article}
\usepackage[utf8]{inputenc}
\usepackage[spanish]{babel}
\usepackage{geometry}
\usepackage{graphicx}
\usepackage{booktabs}
\usepackage{array}
\usepackage{longtable}
\usepackage{caption}
\usepackage{hyperref}
\usepackage{amsmath}
\usepackage{amssymb}
\usepackage{listings}
\usepackage{xcolor}

% Configuración de márgenes
\geometry{top=2.5cm, bottom=2.5cm, left=3cm, right=3cm}

% Configuración de títulos
\usepackage{titlesec}
\titleformat{\section}
  {\normalfont\Large\bfseries}{\thesection}{1em}{}
\titleformat{\subsection}
  {\normalfont\large\bfseries}{\thesubsection}{1em}{}

% Hipervínculos
\hypersetup{
    colorlinks=true,
    linkcolor=blue,
    citecolor=blue,
    urlcolor=blue,
}

% Comando para nombre del autor (evita problemas con corchetes)
\newcommand{\authorname}{Jesús Alberto Mora Lozada}
\newcommand{\institucion}{Universidad de Carabobo}

% Inicio del documento
\begin{document}

% ----------------------------------------------------------------
% CARÁTULA
% ----------------------------------------------------------------
\begin{titlepage}
    \centering
    \vspace*{1cm}
    
    \Huge
    \textbf{Informe Técnico}
    
    \vspace{0.5cm}
    \Large
    \textbf{Diseño y Presupuesto de un Servidor para Análisis de Tráfico y Pentesting}
    
    \vspace{1.5cm}
    \normalsize
    \textbf{Presentado por:} \\
    \vspace{0.3cm}
    \authorname
    
    \vspace{0.8cm}
    \textbf{Fecha:} \\
    \vspace{0.3cm}
    22 de Mayo de 2026
    
    \vspace{1.5cm}
    \textbf{Palabras clave:} Servidor Xeon, Pentesting, Virtualización, Análisis de tráfico, Presupuesto técnico
    
    \vfill
    \rule{0.8\textwidth}{0.4pt}
    \vspace{0.3cm}
    \small
    \textbf{Nota académica:} Este informe ha sido elaborado con fines educativos y de planificación técnica. Los precios corresponden al mercado de componentes reacondicionados de grado empresarial (fecha aproximada: 2025).
    
\end{titlepage}

% ----------------------------------------------------------------
% TABLA DE CONTENIDO
% ----------------------------------------------------------------
\tableofcontents
\newpage

% ----------------------------------------------------------------
% INTRODUCCIÓN
% ----------------------------------------------------------------
\section{Introducción y Objetivo}

El presente informe técnico tiene como objetivo describir el diseño y la justificación de un servidor especializado para dos tareas fundamentales en ciberseguridad ofensiva y defensiva:

\begin{enumerate}
    \item \textbf{Análisis de tráfico de red en tiempo real:} Captura, inspección y reproducción de paquetes para detección de intrusiones, análisis forense y generación de métricas.
    \item \textbf{Pentesting (pruebas de penetración):} Ejecución de máquinas virtuales (VMs) con herramientas como Kali Linux, Metasploit, Hashcat y entornos Active Directory.
\end{enumerate}

El servidor debe soportar múltiples máquinas virtuales de forma concurrente, priorizando la estabilidad, la baja latencia de memoria, el alto número de hilos de ejecución, la velocidad de escritura y la ausencia de cuellos de botella. El presupuesto máximo asignado es de \textbf{3000 USD}.

El diseño se basa en la arquitectura \textbf{Intel Xeon Scalable (Género 2)} utilizando componentes reacondicionados de grado empresarial, que ofrecen la mejor relación rendimiento/precio para este rango.

\section{Parámetros técnicos de diseño}

A continuación se analizan los ocho parámetros solicitados, justificando cada elección de componente.

\subsection{Placa madre con estabilidad amplia}

\begin{itemize}
    \item \textbf{Componente seleccionado:} Dell PowerEdge R740 (placa base propietaria Dell, formato 2U).
    \item \textbf{Características clave:}
        \begin{itemize}
            \item Capacitadores de estado sólido de grado industrial.
            \item Trazados de señal optimizados para reducir diafonía (crosstalk) entre líneas de alta velocidad.
            \item Sistemas de monitoreo de voltaje por componente (sensores integrados en VRM).
            \item Certificación MIL-STD-810G para vibración y choque.
            \item Soporte ECC (Error Correcting Code) en el bus de memoria.
        \end{itemize}
    \item \textbf{Estabilidad:} Diseñada para ciclo de trabajo 24/7/365 con temperatura ambiente de hasta 35°C.
    \item \textbf{Presupuesto asignado:} 950 USD (incluye chasis, backplane, fuentes redundantes y placa).
\end{itemize}

\subsection{Latencia de la memoria y velocidad}

\begin{itemize}
    \item \textbf{Configuración:} 128 GB DDR4 ECC Registered (RDIMM) en disposición \textbf{8 x 16 GB}.
    \item \textbf{Velocidad efectiva:} 2933 MHz (compatible con los Xeon Gold 6130).
    \item \textbf{Latencia CAS típica:} CL17 – CL19.
    \item \textbf{Optimización de canales:}
        \begin{itemize}
            \item Cada CPU Xeon Gold 6130 dispone de \textbf{6 canales de memoria}.
            \item Con 8 módulos por socket (16 totales) se llenan todos los canales, maximizando el ancho de banda y minimizando la latencia de acceso.
        \end{itemize}
    \item \textbf{Impacto en pentesting:} Menor latencia = mejor rendimiento en ataques de diccionario (hashcat) y captura de paquetes con timestamps precisos (Wireshark en modo monitor).
    \item \textbf{Presupuesto asignado:} 350 USD.
\end{itemize}

\subsection{Priorización de hilos y núcleos}

\begin{itemize}
    \item \textbf{Procesadores:} 2 x Intel Xeon Gold 6130.
    \item \textbf{Arquitectura:}
        \[
        \text{16 núcleos} \times 2 = 32 \text{ núcleos físicos}
        \]
        \[
        \text{32 hilos lógicos} \times 2 = 64 \text{ hilos totales}
        \]
    \item \textbf{Frecuencia:} Base 2.1 GHz / Turbo máx. 3.7 GHz.
    \item \textbf{Caché:} 22 MB de caché L3 por CPU.
    \item \textbf{Tecnologías de priorización:}
        \begin{itemize}
            \item Hyper-Threading (2 hilos por núcleo).
            \item Intel VT-x con Extended Page Tables (EPT) para virtualización eficiente.
            \item Intel Cache Allocation Technology (CAT) – permite reservar porciones de caché L3 para VMs críticas.
        \end{itemize}
    \item \textbf{Asignación recomendada en hipervisor:} Fijar núcleos dedicados a VMs de análisis de tráfico (Zeek, Snort) y pentesting (Kali con ataques de fuerza bruta).
    \item \textbf{Presupuesto asignado:} 400 USD (ambos CPUs, mercado refurbished).
\end{itemize}

\subsection{Velocidad de escritura (Storage I/O)}

\begin{table}[!ht]
    \centering
    \caption{Sistema de almacenamiento jerarquizado}
    \label{tab:storage}
    \begin{tabular}{p{4cm}p{6cm}p{3cm}}
        \toprule
        \textbf{Nivel} & \textbf{Configuración} & \textbf{Velocidad estimada} \\
        \midrule
        SO y metadatos & RAID 1: 2 x SSD SATA Enterprise 480 GB & Lectura: 550 MB/s \\
        VMs activas & RAID 10: 4 x SSD SATA Enterprise 960 GB & Escritura secuencial: $>$2000 MB/s \\
        Almacenamiento frío/logs & RAID 1: 2 x HDD SAS 2.4 TB 10k RPM & IOPS: $\sim$300k aleatorios \\
        \bottomrule
    \end{tabular}
\end{table}

\begin{itemize}
    \item \textbf{Controladora RAID:} Dell PERC H740P (caché de 8 GB con protección por batería, escritura diferida segura).
    \item \textbf{Relevancia en pentesting:}
        \begin{itemize}
            \item Capturas pcap masivas (archivos de varios GB) requieren escritura sostenida sin bloqueos.
            \item Reproducción de tráfico con \texttt{tcpreplay} demanda baja latencia de I/O.
            \item Análisis forense sobre imágenes de disco (VMs) requiere acceso aleatorio rápido.
        \end{itemize}
    \item \textbf{Presupuesto asignado:} 650 USD.
\end{itemize}

\subsection{Fuente de poder – soporte eléctrico}

\begin{itemize}
    \item \textbf{Configuración:} 2 fuentes de 750W Platinum (redundantes, hot-swap).
    \item \textbf{Certificación:} 80 PLUS Platinum (eficiencia $>$92\% al 50\% de carga).
    \item \textbf{Protecciones incluidas:}
        \begin{itemize}
            \item OVP (Over Voltage Protection)
            \item OCP (Over Current Protection)
            \item SCP (Short Circuit Protection)
            \item OTP (Over Temperature Protection)
        \end{itemize}
    \item \textbf{Consumo estimado del sistema:}
        \begin{itemize}
            \item Reposo: $\sim$150W
            \item Carga media (10 VMs ligeras): $\sim$350W – 450W
            \item Pico máximo (todas las VMs a máxima CPU): $\sim$600W
        \end{itemize}
    \item \textbf{Redundancia:} Si falla una fuente, la segunda asume la carga completa sin interrupción – crítico para laboratorios de pentesting de larga duración (24h+).
    \item \textbf{Presupuesto:} Incluido en el chasis (950 USD total).
\end{itemize}

\subsection{Evitar cuello de botella entre componentes}

\begin{table}[!ht]
    \centering
    \caption{Análisis de cuellos de botella potenciales}
    \label{tab:bottleneck}
    \begin{tabular}{p{4.5cm}p{4cm}p{4cm}p{2.5cm}}
        \toprule
        \textbf{Ruta de datos} & \textbf{Ancho de banda teórico} & \textbf{Demanda real estimada} & \textbf{¿Cuello?} \\
        \midrule
        CPU a RAM & 6 canales @ 2933 MHz = $\sim$140 GB/s por CPU & $\sim$50 GB/s (máx) & No \\
        CPU 1 a CPU 2 & 2 enlaces UPI @ 10.4 GT/s & Variable según afinidad & Condicional* \\
        RAID a PCIe & PCIe 3.0 x8 = 7.88 GB/s & $<$ 2 GB/s (SSD SATA) & No \\
        NIC a RAM & 40 Gbps = 5 GB/s & $<$ 4 GB/s & No \\
        GPU a CPU & PCIe 3.0 x16 = 64 GB/s & $\sim$2 GB/s (T400) & No \\
        \bottomrule
    \end{tabular}
    \caption*{*Condicional: Asignar VMs a núcleos del mismo socket para evitar tráfico entre CPUs.}
\end{table}

\begin{itemize}
    \item \textbf{Recomendaciones para evitar cuello:}
        \begin{enumerate}
            \item Usar \textbf{affinity} (fijación de núcleos) en el hipervisor: cada VM pesada debe asignarse a un solo socket.
            \item Evitar mezclar tráfico de red intensivo (10GbE) con acceso masivo a disco en el mismo núcleo.
            \item La GPU (NVIDIA T400) al ser de perfil bajo (x8) no satura el bus x16.
        \end{enumerate}
    \item \textbf{Conclusión:} El único cuello potencial sería el acceso a SSD SATA en lugar de NVMe, pero se mantiene dentro del presupuesto.
\end{itemize}

\subsection{El chip de la placa (Chipset)}

\begin{itemize}
    \item \textbf{Chipset:} Intel C622 (nombre en código Lewisburg).
    \item \textbf{Características técnicas destacadas:}
        \begin{itemize}
            \item \textbf{Líneas PCIe adicionales:} 24 líneas PCIe 3.0 (para dispositivos de gestión como iDRAC, USB, SATA).
            \item \textbf{Intel VMD:} Volume Management Device – permite hot-swap de unidades NVMe sin reiniciar.
            \item \textbf{Intel SGX:} Software Guard Extensions – permite crear enclaves seguros para aislar VMs de análisis malicioso.
            \item \textbf{Gestión avanzada de energía:} Soportes C-states profundos y P-states dinámicos.
        \end{itemize}
    \item \textbf{Impacto práctico:}
        \begin{itemize}
            \item Permite pasar GPUs o NICs a VMs mediante VT-d (Direct I/O) sin intervención del chipset, reduciendo latencia.
            \item El iDRAC (Dell Remote Access Controller) integrado permite gestión fuera de banda, crítica para laboratorios remotos.
        \end{itemize}
    \item \textbf{Limitación:} No soporta PCIe 4.0 (disponible en Xeon Scalable 3ra generación, fuera de presupuesto).
\end{itemize}

\subsection{El socket del procesador}

\begin{itemize}
    \item \textbf{Socket:} LGA 3647 (específicamente P0 para Xeon Gold 6130).
    \item \textbf{Características físicas:}
        \begin{itemize}
            \item \textbf{Número de contactos:} 3647 pines (mayor densidad que LGA 2066).
            \item \textbf{Soporte TDP:} Hasta 205W por CPU (los Gold 6130 son de 125W, amplio margen).
            \item \textbf{Topología:} Dos zócalos por placa base con interconexión mediante bus en malla (mesh).
        \end{itemize}
    \item \textbf{Ventajas técnicas:}
        \begin{enumerate}
            \item Permite actualización futura a Xeon Platinum (más núcleos y caché).
            \item Mayor distancia entre contactos que sockets de escritorio $\rightarrow$ menor riesgo de corrosión o puentes accidentales.
            \item Compatibilidad con memoria Intel Optane Persistent Memory (aunque no se utilice ahora).
        \end{enumerate}
    \item \textbf{Presupuesto:} Incluido en la placa base.
\end{itemize}

\section{Componente adicional: Tarjeta de red y GPU}

Para cumplir con los requisitos de análisis de tráfico y visualización, se incluyen dos componentes adicionales:

\subsection{Tarjeta de red de alta velocidad}

\begin{itemize}
    \item \textbf{Modelo:} Intel X710-DA4.
    \item \textbf{Especificaciones:} 4 puertos SFP+ a 10 GbE cada uno (40 Gbps agregados).
    \item \textbf{Tecnologías clave:}
        \begin{itemize}
            \item SR-IOV (Single Root I/O Virtualization): permite pasar puertos virtuales directamente a VMs.
            \item Soporte para VXLAN, NVGRE y checksum offloading.
        \end{itemize}
    \item \textbf{Presupuesto:} 300 USD (refurbished).
\end{itemize}

\subsection{Tarjeta gráfica básica (salida de video)}

\begin{itemize}
    \item \textbf{Modelo:} NVIDIA Quadro P400 (refurbished).
    \item \textbf{Razón de inclusión:} Los Xeon Gold no tienen gráficos integrados; se necesita una GPU mínima para la instalación del SO.
    \item \textbf{Especificaciones:} 256 núcleos CUDA, 2 GB GDDR5, consumo $<$30W, perfil de media altura.
    \item \textbf{Nota:} No acelera significativamente hashcat (aunque es compatible con CUDA), su función principal es la salida de video.
    \item \textbf{Presupuesto:} 50 USD.
\end{itemize}

\section{Resumen de presupuesto final}

\begin{table}[!ht]
    \centering
    \caption{Presupuesto detallado del servidor completo}
    \label{tab:budget}
    \begin{tabular}{p{5cm}p{6cm}r}
        \toprule
        \textbf{Componente / Concepto} & \textbf{Detalle} & \textbf{Costo (USD)} \\
        \midrule
        Servidor base & Dell PowerEdge R740 (chasis + placa + 2 fuentes 750W Platinum + iDRAC) & 950 \\
        Procesadores & 2 x Intel Xeon Gold 6130 (refurbished) & 400 \\
        Memoria RAM & 128 GB DDR4 ECC RDIMM (8x16 GB, 2933 MHz) & 350 \\
        Almacenamiento & 2 SSD 480 GB + 4 SSD 960 GB + 2 HDD 2.4 TB + controladora PERC H740P & 650 \\
        Tarjeta de red & Intel X710-DA4 (4 puertos 10GbE SFP+) & 300 \\
        GPU básica & NVIDIA Quadro P400 (refurbished) & 50 \\
        \midrule
        \textbf{Subtotal} & & \textbf{2.700} \\
        Margen de contingencia & Cables SFF-8643, pasta térmica, adaptadores & 300 \\
        \midrule
        \textbf{Total máximo} & & \textbf{3.000} \\
        \bottomrule
    \end{tabular}
\end{table}

\section{Consideraciones finales y limitaciones}

\subsection{Ventajas del diseño propuesto}
\begin{itemize}
    \item Excelente relación núcleos/precio (32 núcleos físicos por 3000 USD).
    \item Memoria ECC con latencia controlada, crítica para virtualización y análisis de tráfico.
    \item Redundancia de fuentes y RAID10 para tolerancia a fallos.
    \item Posibilidad de pasar dispositivos PCIe (NIC, GPU) directamente a VMs mediante SR-IOV/VT-d.
\end{itemize}

\subsection{Limitaciones importantes}
\begin{itemize}
    \item \textbf{Ruido:} Los ventiladores del R740 emiten $\sim$40 dB en reposo y hasta $\sim$70 dB en carga. No es apto para un entorno de oficina silencioso o doméstico.
    \item \textbf{PCIe 3.0:} El chipset C622 no soporta PCIe 4.0 (la velocidad máxima es 8 GT/s frente a los 16 GT/s de las generaciones más recientes).
    \item \textbf{Consumo eléctrico:} Aunque eficiente (Platinum), su factura eléctrica será notable si funciona 24/7.
    \item \textbf{Portabilidad:} Peso superior a 25 kg, no es transportable sin ayuda.
\end{itemize}

\subsection{Recomendación académica final}

El diseño presentado demuestra un equilibrio óptimo entre los ocho parámetros técnicos solicitados (estabilidad, latencia de memoria, priorización de hilos, velocidad de escritura, soporte eléctrico, ausencia de cuellos de botella, chipset y socket) respetando estrictamente el presupuesto de 3000 USD mediante el uso inteligente de componentes reacondicionados de grado empresarial.

Se recomienda su implementación en un laboratorio universitario o entorno de investigación con infraestructura de refrigeración y aislamiento acústico adecuados.

\vspace{1cm}

\begin{flushright}
\textit{Firmado:} \\
\authorname \\
\institucion \\
22 de Mayo de 2026
\end{flushright}

\newpage
\section*{Referencias técnicas}

\begin{thebibliography}{99}
    \bibitem{intel_xeon} Intel Corporation. (2017). \textit{Intel Xeon Scalable Processor Family Datasheet, Volume 1}. Document Number: 336212-003US.
    
    \bibitem{dell_r740} Dell EMC. (2018). \textit{Dell EMC PowerEdge R740 Technical Guide}. Dell Doc. No. 4.5.
    
    \bibitem{perf_analysis} Molka, D., et al. (2015). ``Memory performance and cache coherency effects on Intel Xeon systems''. \textit{Proceedings of the 2015 International Conference on Parallel Processing}.
    
    \bibitem{proxmox_passthrough} Proxmox Server Solutions. (2024). \textit{PCI(e) Passthrough - Proxmox VE Documentation}. Recuperado de \url{https://pve.proxmox.com/wiki/PCI(e)_Passthrough}
    
    \bibitem{spec_memory} JEDEC Solid State Technology Association. (2018). \textit{DDR4 SDRAM Standard JESD79-4B}.
\end{thebibliography}

\end{document}
